\chapter{Structures}

Atomic data cannot represent compound information such as a point $(x,y)$, a
student record, or a date. \textbf{Structures} bundle a fixed number of values
into one value.

\section{Defining a structure}

\begin{lstlisting}[style=racket]
(define-struct posn (x y))
\end{lstlisting}
(In fact \code{posn} is built into the teaching languages, but this is how it
would be defined.) A single \code{define-struct} creates several functions:
\begin{itemize}
  \item \textbf{constructor} \code{make-posn}: \code{Num Num -> Posn};
  \item \textbf{selectors} \code{posn-x} and \code{posn-y}: \code{Posn -> Num};
  \item \textbf{type predicate} \code{posn?}: \code{Any -> Bool}.
\end{itemize}

\begin{lstlisting}[style=racket]
(define p (make-posn 3 4))
(posn-x p)        ; => 3
(posn-y p)        ; => 4
(posn? p)         ; => true
(posn? 7)         ; => false
\end{lstlisting}

A value such as \code{(make-posn 3 4)} is itself a \emph{value}: it cannot be
simplified further. Structures are immutable in CS 135; to ``change'' a field
you construct a new structure.

\subsection{Substitution rules for structures}
\begin{important}[Structure rules]
For \code{(define-struct sname (f1 ... fn))}:
\begin{itemize}
  \item \code{(sname-fi (make-sname v1 ... vn))} $\Rightarrow$ \code{vi}
  \item \code{(sname? (make-sname v1 ... vn))} $\Rightarrow$ \code{true}
  \item \code{(sname? v)} $\Rightarrow$ \code{false} for any other value \code{v}
\end{itemize}
\end{important}

\section{Data definitions and structure definitions}

Every \code{define-struct} is accompanied by a comment called the
\textbf{structure definition} (field names and types) and a \textbf{data
definition} giving the type a capitalized name:

\begin{lstlisting}[style=racket]
(define-struct student (name id avg))
;; A Student is a (make-student Str Nat Num)
;;   requires: 0 <= avg <= 100
\end{lstlisting}

\section{Templates for structures}

A template for a function consuming a structure selects each field, since a
function consuming a Student will usually need some subset of its fields:

\begin{lstlisting}[style=racket]
;; student-template: Student -> Any
(define (student-template s)
  (... (student-name s) ...
       (student-id s) ...
       (student-avg s) ...))
\end{lstlisting}

\begin{example}[Distance between points]
\begin{lstlisting}[style=racket]
;; (distance p1 p2) produces the Euclidean distance between p1 and p2.
;; Examples:
(check-expect (distance (make-posn 0 0) (make-posn 3 4)) 5)
(check-expect (distance (make-posn 1 1) (make-posn 1 1)) 0)

;; distance: Posn Posn -> Num
(define (distance p1 p2)
  (sqrt (+ (sqr (- (posn-x p1) (posn-x p2)))
           (sqr (- (posn-y p1) (posn-y p2))))))

;; Tests:
(check-within (distance (make-posn 0 0) (make-posn 1 1)) (sqrt 2) 0.0001)
\end{lstlisting}
\end{example}

\begin{example}[Producing a new structure]
\begin{lstlisting}[style=racket]
;; (curve-avg s bonus) produces a copy of s with bonus added to the
;;   average, capped at 100.
;; curve-avg: Student Num -> Student
(check-expect (curve-avg (make-student "Ada" 1 70) 5)
              (make-student "Ada" 1 75))
(check-expect (curve-avg (make-student "Bo" 2 98) 5)
              (make-student "Bo" 2 100))
(define (curve-avg s bonus)
  (make-student (student-name s)
                (student-id s)
                (min 100 (+ (student-avg s) bonus))))
\end{lstlisting}
\end{example}

\paragraph{Trace.}
\begin{lstlisting}[style=racket]
(curve-avg (make-student "Ada" 1 70) 5)
=> (make-student (student-name (make-student "Ada" 1 70))
                 (student-id (make-student "Ada" 1 70))
                 (min 100 (+ (student-avg (make-student "Ada" 1 70)) 5)))
=> (make-student "Ada"
                 (student-id (make-student "Ada" 1 70))
                 (min 100 (+ (student-avg (make-student "Ada" 1 70)) 5)))
=> (make-student "Ada" 1 (min 100 (+ (student-avg (make-student "Ada" 1 70)) 5)))
=> (make-student "Ada" 1 (min 100 (+ 70 5)))
=> (make-student "Ada" 1 (min 100 75))
=> (make-student "Ada" 1 75)
\end{lstlisting}

\section{Nested structures}

Fields can themselves be structures. The template then calls the template of
the inner type (or a helper) on that field.

\begin{lstlisting}[style=racket]
(define-struct circle (centre radius))
;; A Circle is a (make-circle Posn Num)
;;   requires: radius > 0

;; circle-template: Circle -> Any
(define (circle-template c)
  (... (posn-template (circle-centre c)) ...
       (circle-radius c) ...))

;; (inside? c p) produces true if p lies inside or on circle c.
;; inside?: Circle Posn -> Bool
(check-expect (inside? (make-circle (make-posn 0 0) 5) (make-posn 3 4)) true)
(check-expect (inside? (make-circle (make-posn 0 0) 5) (make-posn 4 4)) false)
(define (inside? c p)
  (<= (distance (circle-centre c) p) (circle-radius c)))
\end{lstlisting}

\section{Mixed data with structures}

A type may be one of several structure types. The data definition uses
\code{anyof} and the template uses the type predicates in a \code{cond}:

\begin{lstlisting}[style=racket]
(define-struct rect (corner width height))
;; A Rect is a (make-rect Posn Num Num)
;; A Shape is (anyof Circle Rect)

;; shape-template: Shape -> Any
(define (shape-template sh)
  (cond [(circle? sh) (... (circle-centre sh) ... (circle-radius sh) ...)]
        [(rect? sh) (... (rect-corner sh) ... (rect-width sh) ...
                         (rect-height sh) ...)]))

;; (area sh) produces the area of shape sh.
;; area: Shape -> Num
(check-expect (area (make-rect (make-posn 0 0) 2 3)) 6)
(define (area sh)
  (cond [(circle? sh) (* pi (sqr (circle-radius sh)))]
        [(rect? sh) (* (rect-width sh) (rect-height sh))]))
\end{lstlisting}

\begin{aside}
Racket is \emph{dynamically typed}: nothing stops you from writing
\code{(make-posn "a" 'b)}. Contracts and data definitions are comments that
document intended use. Selectors do check their argument's structure type,
so \code{(posn-x 5)} is a run-time error.
\end{aside}

\section{Common pitfalls}
\begin{itemize}
  \item Forgetting the data definition comment or field types.
  \item Writing \code{(x p)} instead of \code{(posn-x p)}.
  \item Trying to mutate a field; construct a new structure instead.
  \item Comparing structures with \code{=}; use \code{equal?} (and \code{check-expect} compares structurally).
\end{itemize}

\section{Summary}
\begin{itemize}
  \item \code{define-struct} creates a constructor, selectors and a predicate.
  \item Every structure has a structure definition and a data definition comment.
  \item The structure template selects every field; nested structures lead to nested template calls.
  \item Mixed data uses \code{anyof} and type predicates in a \code{cond}.
\end{itemize}
